% =========================================================
% Limits at Infinity --- Notes
% Chapter: continuity
% Volume III - Analysis
% =========================================================

\subsection*{Limits at Infinity}
\label{sec:limits-at-infinity}
\begin{definitionbox}{Definition (Infinite Adherent Points)}
\begin{definition}[Infinite Adherent Points]
\label{def:infinite-adherent-points}
Let $X\subseteq\mathbb{R}$. The point $+\infty$ is adherent to $X$ if
and only if for every $M\in\mathbb{R}$ there exists $x\in X$ such that
$x>M$. The point $-\infty$ is adherent to $X$ if and only if for every
$M\in\mathbb{R}$ there exists $x\in X$ such that $x<M$.
\end{definition}
\end{definitionbox}
\LeanFormalizes{def:infinite-adherent-points}{lra-lean}{LRA.VolumeIII.Analysis.Continuity.LimitsAtInfinityBase}{PlusInftyAdherent}{pending}

\begin{remark*}[Standard quantified statement]
\[
\begin{aligned}
&+\infty\text{ is adherent to }X
\Longleftrightarrow
\forall M\in\mathbb{R}\;\exists x\in X\;(x>M),\\
&-\infty\text{ is adherent to }X
\Longleftrightarrow
\forall M\in\mathbb{R}\;\exists x\in X\;(x<M).
\end{aligned}
\]
\end{remark*}

\begin{remark*}[Predicate reading]
\[
\begin{aligned}
\left(+\infty\in\operatorname{Closure}_{\overline{\mathbb{R}}}(X)\right)
&\coloneqq
\forall M\in\mathbb{R}\;\exists x\in X\;(x>M),\\
\left(-\infty\in\operatorname{Closure}_{\overline{\mathbb{R}}}(X)\right)
&\coloneqq
\forall M\in\mathbb{R}\;\exists x\in X\;(x<M).
\end{aligned}
\]
\end{remark*}

\begin{remark*}[Negated quantified statement]
\[
\begin{aligned}
&\neg(+\infty\text{ is adherent to }X)
\Longleftrightarrow
\exists M\in\mathbb{R}\;\forall x\in X\;(x\le M),\\
&\neg(-\infty\text{ is adherent to }X)
\Longleftrightarrow
\exists M\in\mathbb{R}\;\forall x\in X\;(x\ge M).
\end{aligned}
\]
\end{remark*}

\begin{remark*}[Negation predicate reading]
\[
\begin{aligned}
\neg\left(+\infty\in\operatorname{Closure}_{\overline{\mathbb{R}}}(X)\right)
&\Longleftrightarrow
\operatorname{BoundedAbove}(X,\mathbb{R}),\\
\neg\left(-\infty\in\operatorname{Closure}_{\overline{\mathbb{R}}}(X)\right)
&\Longleftrightarrow
\operatorname{BoundedBelow}(X,\mathbb{R}).
\end{aligned}
\]
\end{remark*}

\begin{remark*}[Failure modes]
\begin{description}
  \item[Exposition.]
Infinite adherence fails at $+\infty$ exactly when $X$ is bounded above,
and fails at $-\infty$ exactly when $X$ is bounded below.


  \item[\(\operatorname{Closure}.\)]
  The displayed obstruction is the corresponding failure mode.
  \[
\begin{aligned}
&\neg(+\infty\text{ is adherent to }X)
\Longleftrightarrow
\exists M\in\mathbb{R}\;\forall x\in X\;(x\le M),\\
&\neg(-\infty\text{ is adherent to }X)
\Longleftrightarrow
\exists M\in\mathbb{R}\;\forall x\in X\;(x\ge M).
\end{aligned}
  \]
  \[
\begin{aligned}
\neg\left(+\infty\in\operatorname{Closure}_{\overline{\mathbb{R}}}(X)\right)
&\Longleftrightarrow
\operatorname{BoundedAbove}(X,\mathbb{R}),\\
\neg\left(-\infty\in\operatorname{Closure}_{\overline{\mathbb{R}}}(X)\right)
&\Longleftrightarrow
\operatorname{BoundedBelow}(X,\mathbb{R}).
\end{aligned}
  \]
\end{description}
\end{remark*}

\begin{remark*}[Failure modes]
\begin{description}
  \item[Exposition.]
\[
\underbrace{\exists M\in\mathbb{R}\;\forall x\in X\;(x\le M)}_{\text{finite upper bound}}
\qquad
\underbrace{\exists M\in\mathbb{R}\;\forall x\in X\;(x\ge M)}_{\text{finite lower bound}}.
\]


  \item[\(\operatorname{Closure}.\)]
  The displayed obstruction is the corresponding failure mode.
  \[
\begin{aligned}
&\neg(+\infty\text{ is adherent to }X)
\Longleftrightarrow
\exists M\in\mathbb{R}\;\forall x\in X\;(x\le M),\\
&\neg(-\infty\text{ is adherent to }X)
\Longleftrightarrow
\exists M\in\mathbb{R}\;\forall x\in X\;(x\ge M).
\end{aligned}
  \]
  \[
\begin{aligned}
\neg\left(+\infty\in\operatorname{Closure}_{\overline{\mathbb{R}}}(X)\right)
&\Longleftrightarrow
\operatorname{BoundedAbove}(X,\mathbb{R}),\\
\neg\left(-\infty\in\operatorname{Closure}_{\overline{\mathbb{R}}}(X)\right)
&\Longleftrightarrow
\operatorname{BoundedBelow}(X,\mathbb{R}).
\end{aligned}
  \]
\end{description}
\end{remark*}

\begin{remark*}[Interpretation]
Both finite and infinite adherence are instances of adherence in
$\overline{\mathbb{R}}$: finite adherence uses balls; infinite
adherence uses rays.
\end{remark*}

\begin{dependencies}
\begin{itemize}
  \item \hyperref[def:extended-real-numbers]{Extended Real Numbers}
  \item \hyperref[def:reals]{Real Numbers}
  \item \hyperref[def:subset]{Subset}
  \item \hyperref[def:ray]{Ray}
\end{itemize}
\end{dependencies}
\begin{definitionbox}{Definition (Limit at Infinity)}
\begin{definition}[Limit at Infinity]
\label{def:limit-at-infinity}
Let $X \subseteq \mathbb{R}$ be such that for every $M \in \mathbb{R}$ there exists $x \in X$ with $x > M$, let $f : X \to \mathbb{R}$, and let $L \in \mathbb{R}$. The limit of $f$ at $+\infty$ equals $L$ if and only if
for every $\varepsilon > 0$ there exists $M \in \mathbb{R}$ such that for every $x \in X$, $x > M$ implies $|f(x) - L| < \varepsilon$.
\end{definition}
\end{definitionbox}
\LeanFormalizes{def:limit-at-infinity}{lra-lean}{LRA.VolumeIII.Analysis.Continuity.LimitsAtInfinityAdditions}{TendsToInfty}{pending}
\LeanFormalizes{def:limit-at-infinity}{lra-lean}{LRA.VolumeIII.Analysis.Continuity.LimitsAtInfinityAdditions}{TendsToNegInfty}{pending}

\begin{remark*}[Standard quantified statement]
\[
\begin{aligned}
&\text{Let } X \subseteq \mathbb{R},\ \forall M \in \mathbb{R}\ \exists x \in X\ (x > M),\ f : X \to \mathbb{R},\ L \in \mathbb{R}.\\
&\lim_{x \to +\infty} f(x) = L \Longleftrightarrow \forall \varepsilon > 0\ \exists M \in \mathbb{R}\ \forall x \in X\\
&\qquad (x > M \Rightarrow |f(x) - L| < \varepsilon).
\end{aligned}
\]
\end{remark*}

\begin{remark*}[Predicate reading]
\[
\operatorname{LimitAtInfinity}(f,+\infty,L) \coloneqq \forall \varepsilon > 0\ \exists M \in \mathbb{R}\ \forall x \in X\ (x > M \Rightarrow |f(x) - L| < \varepsilon).
\]
\end{remark*}

\begin{remark*}[Negated quantified statement]
\[
\begin{aligned}
&\text{Let } X \subseteq \mathbb{R},\ \forall M \in \mathbb{R}\ \exists x \in X\ (x > M),\ f : X \to \mathbb{R},\ L \in \mathbb{R}.\\
&\neg\bigl(\lim_{x \to +\infty} f(x) = L\bigr) \Longleftrightarrow \exists \varepsilon_0 > 0\ \forall M \in \mathbb{R}\ \exists x \in X\\
&\qquad (x > M \land |f(x) - L| \ge \varepsilon_0).
\end{aligned}
\]
\end{remark*}

\begin{remark*}[Negation predicate reading]
\[
\neg\operatorname{LimitAtInfinity}(f,+\infty,L) \Longleftrightarrow \exists \varepsilon_0 > 0\ \forall M \in \mathbb{R}\ \exists x \in X\ (x > M \land |f(x) - L| \ge \varepsilon_0).
\]
\end{remark*}

\begin{remark*}[Failure modes]
\begin{description}
  \item[Exposition.]
The definition fails only when there is a fixed tolerance $\varepsilon_0$ that the function cannot satisfy on any tail of $X$, meaning every attempt to choose $M$ leaves some point $x > M$ with $|f(x) - L| \ge \varepsilon_0$.


  \item[\(\operatorname{LimitAtInfinity}.\)]
  The displayed obstruction is the corresponding failure mode.
  \[
\begin{aligned}
&\text{Let } X \subseteq \mathbb{R},\ \forall M \in \mathbb{R}\ \exists x \in X\ (x > M),\ f : X \to \mathbb{R},\ L \in \mathbb{R}.\\
&\neg\bigl(\lim_{x \to +\infty} f(x) = L\bigr) \Longleftrightarrow \exists \varepsilon_0 > 0\ \forall M \in \mathbb{R}\ \exists x \in X\\
&\qquad (x > M \land |f(x) - L| \ge \varepsilon_0).
\end{aligned}
  \]
  \[
\neg\operatorname{LimitAtInfinity}(f,+\infty,L) \Longleftrightarrow \exists \varepsilon_0 > 0\ \forall M \in \mathbb{R}\ \exists x \in X\ (x > M \land |f(x) - L| \ge \varepsilon_0).
  \]
\end{description}
\end{remark*}

\begin{remark*}[Failure modes]
\begin{description}
  \item[Exposition.]
\[
\underbrace{\exists \varepsilon_0 > 0}_{\text{fixed tolerance}}
\underbrace{\forall M \in \mathbb{R}}_{\text{no tail succeeds}}
\underbrace{\exists x \in X\ (x > M \land |f(x) - L| \ge \varepsilon_0)}_{\text{far-point violation}}.
\]


  \item[\(\operatorname{LimitAtInfinity}.\)]
  The displayed obstruction is the corresponding failure mode.
  \[
\begin{aligned}
&\text{Let } X \subseteq \mathbb{R},\ \forall M \in \mathbb{R}\ \exists x \in X\ (x > M),\ f : X \to \mathbb{R},\ L \in \mathbb{R}.\\
&\neg\bigl(\lim_{x \to +\infty} f(x) = L\bigr) \Longleftrightarrow \exists \varepsilon_0 > 0\ \forall M \in \mathbb{R}\ \exists x \in X\\
&\qquad (x > M \land |f(x) - L| \ge \varepsilon_0).
\end{aligned}
  \]
  \[
\neg\operatorname{LimitAtInfinity}(f,+\infty,L) \Longleftrightarrow \exists \varepsilon_0 > 0\ \forall M \in \mathbb{R}\ \exists x \in X\ (x > M \land |f(x) - L| \ge \varepsilon_0).
  \]
\end{description}
\end{remark*}

\begin{remark*}[Interpretation]
The limit at $+\infty$ demands that eventually the function values stay within any prescribed accuracy of $L$. Intuitively, once $x$ is large enough, the graph of $f$ must lie in the horizontal band of height $\varepsilon$ around $L$. Failure occurs when some positive tolerance is violated infinitely often by points escaping to $+\infty$. This definition treats unbounded domains using the same epsilon control that governs finite limits, capturing the long-range behavior of $f$ with respect to the ideal point at infinity.
\end{remark*}

\begin{dependencies}
\hyperref[def:infinite-adherent-points]{Definition~\ref*{def:infinite-adherent-points}}.
\end{dependencies}

% Lean-aligned additions inserted for Volume III Lean-to-TeX handoff.
% Lean-aligned addition: thm:limit-at-neg-infinity-iff-reflection
\begin{theorem}[Limit at Negative Infinity by Reflection]
\label{thm:limit-at-neg-infinity-iff-reflection}
A function $f$ has limit $L$ as $x\to -\infty$ if and only if the reflected function $x\mapsto f(-x)$ has limit $L$ as $x\to +\infty$.

\smallskip
\noindent
\hyperref[prf:limit-at-neg-infinity-iff-reflection]{\textit{Go to proof.}}
\end{theorem}
\LeanFormalizes{thm:limit-at-neg-infinity-iff-reflection}{lra-lean}{LRA.VolumeIII.Analysis.Continuity.LimitsAtInfinityAdditions}{LimitAtNegInfinityIffReflection}{pending}

\begin{remark*}[Standard quantified statement]
A function $f$ has limit $L$ as $x\to -\infty$ if and only if the reflected function $x\mapsto f(-x)$ has limit $L$ as $x\to +\infty$.
\end{remark*}

\begin{remark*}[Interpretation]
This formal object records a Lean-aligned addition in the local language of this section.
\end{remark*}

\begin{dependencies}
\NoLocalDependencies
\end{dependencies}
% Lean-aligned addition: thm:tendsto-infty-add
\begin{theorem}[Sum Law for Limits at Positive Infinity]
\label{thm:tendsto-infty-add}
If $f(x)\to L$ and $g(x)\to M$ as $x\to+\infty$, then $(f+g)(x)\to L+M$ as $x\to+\infty$.

\smallskip
\noindent
\hyperref[prf:tendsto-infty-add]{\textit{Go to proof.}}
\end{theorem}
\LeanFormalizes{thm:tendsto-infty-add}{lra-lean}{LRA.VolumeIII.Analysis.Continuity.LimitsAtInfinityAdditions}{TendstoInftyAdd}{pending}

\begin{remark*}[Standard quantified statement]
If $f(x)\to L$ and $g(x)\to M$ as $x\to+\infty$, then $(f+g)(x)\to L+M$ as $x\to+\infty$.
\end{remark*}

\begin{remark*}[Interpretation]
This formal object records a Lean-aligned addition in the local language of this section.
\end{remark*}

\begin{dependencies}
\NoLocalDependencies
\end{dependencies}
% Lean-aligned addition: thm:tendsto-infty-sub
\begin{theorem}[Difference Law for Limits at Positive Infinity]
\label{thm:tendsto-infty-sub}
If $f(x)\to L$ and $g(x)\to M$ as $x\to+\infty$, then $(f-g)(x)\to L-M$ as $x\to+\infty$.

\smallskip
\noindent
\hyperref[prf:tendsto-infty-sub]{\textit{Go to proof.}}
\end{theorem}
\LeanFormalizes{thm:tendsto-infty-sub}{lra-lean}{LRA.VolumeIII.Analysis.Continuity.LimitsAtInfinityAdditions}{TendstoInftySub}{pending}

\begin{remark*}[Standard quantified statement]
If $f(x)\to L$ and $g(x)\to M$ as $x\to+\infty$, then $(f-g)(x)\to L-M$ as $x\to+\infty$.
\end{remark*}

\begin{remark*}[Interpretation]
This formal object records a Lean-aligned addition in the local language of this section.
\end{remark*}

\begin{dependencies}
\NoLocalDependencies
\end{dependencies}
% Lean-aligned addition: thm:tendsto-infty-scalar
\begin{theorem}[Scalar Law for Limits at Positive Infinity]
\label{thm:tendsto-infty-scalar}
If $f(x)\to L$ as $x\to+\infty$ and $\alpha\in\mathbb R$, then $(\alpha f)(x)\to\alpha L$ as $x\to+\infty$.

\smallskip
\noindent
\hyperref[prf:tendsto-infty-scalar]{\textit{Go to proof.}}
\end{theorem}
\LeanFormalizes{thm:tendsto-infty-scalar}{lra-lean}{LRA.VolumeIII.Analysis.Continuity.LimitsAtInfinityAdditions}{TendstoInftyScalar}{pending}

\begin{remark*}[Standard quantified statement]
If $f(x)\to L$ as $x\to+\infty$ and $\alpha\in\mathbb R$, then $(\alpha f)(x)\to\alpha L$ as $x\to+\infty$.
\end{remark*}

\begin{remark*}[Interpretation]
This formal object records a Lean-aligned addition in the local language of this section.
\end{remark*}

\begin{dependencies}
\NoLocalDependencies
\end{dependencies}
% Lean-aligned addition: thm:tendsto-infty-mul
\begin{theorem}[Product Law for Limits at Positive Infinity]
\label{thm:tendsto-infty-mul}
If $f(x)\to L$ and $g(x)\to M$ as $x\to+\infty$, then $(fg)(x)\to LM$ as $x\to+\infty$.

\smallskip
\noindent
\hyperref[prf:tendsto-infty-mul]{\textit{Go to proof.}}
\end{theorem}
\LeanFormalizes{thm:tendsto-infty-mul}{lra-lean}{LRA.VolumeIII.Analysis.Continuity.LimitsAtInfinityAdditions}{TendstoInftyMul}{pending}

\begin{remark*}[Standard quantified statement]
If $f(x)\to L$ and $g(x)\to M$ as $x\to+\infty$, then $(fg)(x)\to LM$ as $x\to+\infty$.
\end{remark*}

\begin{remark*}[Interpretation]
This formal object records a Lean-aligned addition in the local language of this section.
\end{remark*}

\begin{dependencies}
\NoLocalDependencies
\end{dependencies}
% Lean-aligned addition: thm:tendsto-infty-div
\begin{theorem}[Quotient Law for Limits at Positive Infinity]
\label{thm:tendsto-infty-div}
If $f(x)\to L$ and $g(x)\to M$ as $x\to+\infty$, $M\ne0$, and $g$ is eventually nonzero, then $(f/g)(x)\to L/M$ as $x\to+\infty$.

\smallskip
\noindent
\hyperref[prf:tendsto-infty-div]{\textit{Go to proof.}}
\end{theorem}
\LeanFormalizes{thm:tendsto-infty-div}{lra-lean}{LRA.VolumeIII.Analysis.Continuity.LimitsAtInfinityAdditions}{TendstoInftyDiv}{pending}

\begin{remark*}[Standard quantified statement]
If $f(x)\to L$ and $g(x)\to M$ as $x\to+\infty$, $M\ne0$, and $g$ is eventually nonzero, then $(f/g)(x)\to L/M$ as $x\to+\infty$.
\end{remark*}

\begin{remark*}[Interpretation]
This formal object records a Lean-aligned addition in the local language of this section.
\end{remark*}

\begin{dependencies}
\NoLocalDependencies
\end{dependencies}
% Lean-aligned addition: thm:sequential-criterion-tendsto-infty
\begin{theorem}[Sequential Criterion for Limits at Positive Infinity]
\label{thm:sequential-criterion-tendsto-infty}
A function $f$ has limit $L$ as $x\to+\infty$ if and only if $f(x_n)\to L$ for every sequence $(x_n)$ in the domain that escapes to $+\infty$.

\smallskip
\noindent
\hyperref[prf:sequential-criterion-tendsto-infty]{\textit{Go to proof.}}
\end{theorem}
\LeanFormalizes{thm:sequential-criterion-tendsto-infty}{lra-lean}{LRA.VolumeIII.Analysis.Continuity.LimitsAtInfinityAdditions}{SequentialCriterionTendstoInfty}{pending}

\begin{remark*}[Standard quantified statement]
A function $f$ has limit $L$ as $x\to+\infty$ if and only if $f(x_n)\to L$ for every sequence $(x_n)$ in the domain that escapes to $+\infty$.
\end{remark*}

\begin{remark*}[Interpretation]
This formal object records a Lean-aligned addition in the local language of this section.
\end{remark*}

\begin{dependencies}
\NoLocalDependencies
\end{dependencies}
% Lean-aligned addition: def:escapes-to-infty
\begin{definition}[Sequence Escaping to Positive Infinity]
\label{def:escapes-to-infty}
A real sequence $(x_n)$ \emph{escapes to $+\infty$} if for every $M\in\mathbb R$ there exists $N$ such that $n\ge N$ implies $x_n>M$.
\end{definition}
\LeanFormalizes{def:escapes-to-infty}{lra-lean}{LRA.VolumeIII.Analysis.Continuity.LimitsAtInfinityAdditions}{EscapesToInfty}{pending}

\begin{remark*}[Standard quantified statement]
A real sequence $(x_n)$ \emph{escapes to $+\infty$} if for every $M\in\mathbb R$ there exists $N$ such that $n\ge N$ implies $x_n>M$.
\end{remark*}

\begin{remark*}[Interpretation]
This formal object records a Lean-aligned addition in the local language of this section.
\end{remark*}

\begin{dependencies}
\NoLocalDependencies
\end{dependencies}
% Lean-aligned addition: def:minus-infty-adherent
\begin{definition}[Negative Infinity Adherent Point]
\label{def:minus-infty-adherent}
The ideal point $-\infty$ is adherent to a set $A\subseteq\mathbb R$ if for every $M\in\mathbb R$ there exists $x\in A$ with $x<M$.
\end{definition}
\LeanFormalizes{def:minus-infty-adherent}{lra-lean}{LRA.VolumeIII.Analysis.Continuity.LimitsAtInfinityBase}{MinusInftyAdherent}{pending}

\begin{remark*}[Standard quantified statement]
The ideal point $-\infty$ is adherent to a set $A\subseteq\mathbb R$ if for every $M\in\mathbb R$ there exists $x\in A$ with $x<M$.
\end{remark*}

\begin{remark*}[Interpretation]
This formal object records a Lean-aligned addition in the local language of this section.
\end{remark*}

\begin{dependencies}
\NoLocalDependencies
\end{dependencies}
